\section{Introduction}
\label{sec:intro}

Memory is essential for agents to retain past experiences and use them in future interactions. A personal assistant, for example, must recall what a user said, planned, and shared across months of conversation to answer later questions. Many memory systems therefore decide what to store, how to organize it, and what to retrieve, for instance by scoring a memory stream, moving context between memory tiers, or extracting facts and notes from the dialogue~\citep{park2023generative,packer2023memgpt,chhikara2025mem0,xu2025amem,wang2025mirix}.

Recent studies show that, as language and embedding models become more capable, simple retrieval-augmented generation (RAG) can answer most text-memory questions correctly and outperform most of these systems~\citep{bao2026rawrag,zhou2026explicit,sharma2026selection,peng2026memorytax}. RAG stores the history unchanged, retrieves the passages most similar to the question, and gives them to a reader. We confirm this on LoCoMo~\citep{maharana2024locomo} and LongMemEval-S~\citep{wu2025longmemeval}: with \texttt{gemini-3.8-flash} as the reader, RAG answers about 90\% of the questions correctly and comes within one point of reading the full history (Figure~\ref{fig:rag_motivation}).

However, this level of performance does not consistently carry over to multimodal memory, where users also share receipts, screenshots, and photographs whose content may never appear in the conversation text. With the same reader, RAG reaches only 65 to 81\% accuracy on three multimodal benchmarks and trails full-context reading by 3.9 to 16.1 points (Figure~\ref{fig:rag_motivation}). A larger budget of 40 rounds nearly closes the gap only on MemLens~\citep{ren2026memlens}, whose histories are short, while top-40 retrieval still trails by 4.6 points on MemEye~\citep{guo2026memeye} and 14.1 points on VisualMem~\citep{nguyen2026personalvisual}. The five memory systems we evaluate do not close this gap either (Section~\ref{sec:res-main}).

One main reason is that visual evidence often spans several related images, some of which contain crucial information but have low embedding similarity to the question. To total the expenses of a trip, for example, an assistant must find every relevant receipt, while a fixed number of similar rounds may include only some of them. On MemEye, top-10 retrieval finds all the evidence for 93\% of single-round questions but for only 38\% of questions that require four or five rounds (Figure~\ref{fig:setsize}). Relevance can also depend on context elsewhere in the history: a phrase such as ``the D9 summary'' may identify a card through an earlier exchange rather than through words in the card's own round or image description. Similarity-based retrieval can therefore overlook these images, leaving the reader without the evidence needed to connect relevant facts and derive the answer. Moreover, each image can carry rich information, so missing even one relevant image can leave a substantial gap. On VisualMem, 55\% of the answer-bearing images rank below the 40 rounds most similar to the question, and RAG@40 loses 21 accuracy points on the questions whose evidence it misses. These challenges make reliable multimodal memory difficult to achieve through similarity-based retrieval alone.

To address this problem, we propose Visibility-Gated Set Reading (\vgs{}), which combines \emph{global evidence navigation} with \emph{visibility-guided reading} (Figure~\ref{fig:overview}). Global navigation gives a selector a compact index of the whole history, with one entry per conversation round. The selector reads the entries together, resolves references across rounds, and returns the related rounds as an evidence set whose size it chooses for the question. The reader then opens the original rounds and images, so selection can use context from the entire history while detailed visual information remains available for answering. Because a compact index can omit useful details, the visibility gate measures how much of a history's content the index retains. When visibility is low, it widens evidence access with a reading path suited to the answering model: longer index entries and notes on the selected rounds for a lightweight reader, and wider retrieval for a strong reader. The gate uses one threshold for all benchmarks, and caching and compressing the index reduce the cost of repeated access to a user's memory.

In summary, our contributions are three-fold:
\begin{itemize}
\setlength{\itemsep}{1pt}
\setlength{\parskip}{0pt}
\item We show that the strength of simple RAG on text memory does not carry over to multimodal memory, and trace the gap to image-bearing evidence sets that similarity-based retrieval recovers only in part.
\item We propose \vgs{}, a training-free memory architecture that couples global evidence navigation over a compact history index with visibility-guided reading adapted to the answering model.
\item Experiments on MemEye, MemLens, VisualMem, and CamRoll show that \vgs{} improves a lightweight reader on all four benchmarks, by 4.3 to 12.9 points over RAG, and exceeds re-ranked retrieval, full-context reading, and five recent memory systems. With a strong reader, it is at least as accurate as RAG on all four and gains 9.5 points on VisualMem, where it comes within five points of full-context reading at about a seventh of its cost. Controlled analyses link these gains to more complete evidence sets and to the choice of reading path.
\end{itemize}

\begin{figure*}[t]
\centering
\resizebox{\textwidth}{!}{% Figure 1: overview of Visibility-Gated Set Reading (TikZ, 16 cm wide; included by sections/01_intro.tex)
\begin{tikzpicture}[yscale=1.13,
  font=\fontsize{6.9}{8.2}\selectfont, >={Stealth[length=1.8mm,width=1.5mm]},
  panel/.style={rounded corners=3pt, line width=0.6pt},
  title/.style={font=\fontsize{8.1}{9.5}\selectfont\bfseries, anchor=west, inner sep=0pt},
  num/.style={circle, fill=black!78, text=white, font=\fontsize{5.6}{6}\selectfont\bfseries, inner sep=0pt, minimum size=8pt},
  body/.style={align=left, text width=#1, inner sep=0pt, anchor=north west},
  arr/.style={->, line width=0.8pt, draw=black!70},
  chip/.style={rounded corners=1.2pt, draw=gold!70!black, fill=gold!35, font=\ttfamily\fontsize{5.2}{6}\selectfont, inner sep=1.4pt},
  idx/.style={font=\ttfamily\fontsize{4.6}{5.4}\selectfont, anchor=west, inner sep=0pt, text=black!85},
  note/.style={font=\fontsize{6.0}{7.0}\selectfont\itshape, text=black!62},
]
\colorlet{cMem}{blue!6}\colorlet{cMemD}{blue!50!black}
\colorlet{cGate}{yellow!15}\colorlet{cGateD}{orange!60!black}
\colorlet{cStd}{green!8}\colorlet{cStdD}{green!38!black}
\colorlet{cWid}{orange!10}\colorlet{cWidD}{orange!65!black}
\colorlet{cRead}{violet!7}\colorlet{cReadD}{violet!55!black}
% ---------------------------------------------------------------- phases
\draw[black!40, line width=0.45pt] (0.05,6.02) -- (3.9,6.02);
\node[note, fill=white, inner sep=1.5pt] at (1.97,6.02) {write time};
\draw[black!40, line width=0.45pt] (4.35,6.02) -- (15.95,6.02);
\node[note, fill=white, inner sep=1.5pt] at (10.15,6.02) {question time};
\node[note, text=cStdD] at (9.5,5.62) {navigation over the index};
\node[note, text=cReadD] at (14.05,5.62) {perception of the raw rounds};
% ---------------------------------------------------------------- (1) memory with an index
\draw[panel, draw=cMemD, fill=cMem] (0,0) rectangle (3.95,5.4);
\node[num] at (0.25,5.13) {1}; \node[title] at (0.42,5.13) {Indexed raw memory};
\foreach \y/\t in {4.55/{R37 \textperiodcentered{} 05-06 \textperiodcentered{} dinner at the market}, 4.05/{R12 \textperiodcentered{} 05-03 \textperiodcentered{} booked the hotel}} {
  \draw[rounded corners=1.5pt, draw=black!40, fill=white] (0.22,\y-0.2) rectangle (3.73,\y+0.2);
  \draw[draw=black!40, fill=blue!15] (0.3,\y-0.14) rectangle (0.62,\y+0.14);
  \fill[blue!45] (0.32,\y-0.12) -- (0.43,\y+0.02) -- (0.51,\y-0.05) -- (0.6,\y+0.07) -- (0.6,\y-0.12) -- cycle;
  \fill[yellow!70!orange] (0.55,\y+0.09) circle (0.028);
  \node[idx] at (0.7,\y+0.07) {\t};
  \fill[black!16] (0.7,\y-0.11) rectangle (3.5,\y-0.06);
}
\node[note] at (1.97,3.6) {raw rounds and images, kept unchanged};
\draw[arr, draw=cMemD] (1.97,3.45) -- node[right, note, xshift=1pt] {one line per round} (1.97,3.05);
\draw[rounded corners=2pt, draw=cMemD!75, fill=white, line width=0.5pt] (0.18,0.18) rectangle (3.77,2.95);
\foreach \y/\t in {
  2.70/{R12 05-03 booked the hotel, Alfama},
  2.50/{\hphantom{R12 05-03 }IMG: hotel invoice, 3 nights},
  2.20/{R20 05-04 museums open on Monday?},
  1.90/{R37 05-06 dinner at the market hall},
  1.70/{\hphantom{R37 05-06 }IMG: restaurant bill},
  1.40/{R44 05-07 tram 28 or a walking tour?},
  1.10/{R58 05-09 the D9 summary card},
  0.90/{\hphantom{R58 05-09 }IMG: day summary, receipt},
  0.60/{R61 05-09 taxi to the airport},
  0.40/{\hphantom{R61 05-09 }IMG: taxi receipt}} {
  \node[idx] at (0.42,\y) {\t};
}
\foreach \y in {2.60, 1.80, 1.00, 0.50} { \fill[gold!85!black] (0.3,\y) circle (0.04); }
% ---------------------------------------------------------------- (2) question and visibility gate
\draw[panel, draw=cGateD, fill=cGate] (4.35,0) rectangle (6.95,5.4);
\node[num] at (4.6,5.13) {2}; \node[title] at (4.77,5.13) {Visibility gate};
\node[draw=black!50, fill=white, rounded corners=2.5pt, align=center, text width=2.15cm, inner sep=2.5pt,
      font=\fontsize{6.2}{7.4}\selectfont\itshape] (Q) at (5.65,4.38) {How much did I spend on all receipts from my Lisbon trip?};
\node[draw=cGateD!70, fill=white, rounded corners=2pt, align=left, text width=2.15cm, inner sep=2.5pt] (V) at (5.65,2.85)
      {$H(h)$: the share of each round's content words that its index line shows, averaged over the history; computed at write time};
\node[draw=cGateD, fill=white, diamond, aspect=1.7, inner sep=0.3pt, line width=0.6pt] (D) at (5.65,1.25) {$H(h)\ge\tau$\,?};
\draw[arr] (Q.south) -- (V.north);
\draw[arr] (V.south) -- (D.north);
% ---------------------------------------------------------------- (3) standard path
\draw[panel, draw=cStdD, fill=cStd] (7.35,2.85) rectangle (11.75,5.4);
\node[num] at (7.6,5.13) {3}; \node[title] at (7.77,5.13) {Global evidence navigation};
\node[body=4.1cm] at (7.5,4.85) {A selector reads the index lines of the \emph{whole} history and returns the evidence set $S(q)$, choosing its size for each question.};
\tikzset{gchip/.style={chip, draw=black!35, fill=black!8, text=black!55}}
\node[anchor=west, font=\fontsize{5.8}{6.8}\selectfont] at (7.5,3.48) {cosine top-4};
\node[chip] at (9.0,3.48) {R12}; \node[gchip] at (9.55,3.48) {R20}; \node[gchip] at (10.1,3.48) {R44}; \node[chip] at (10.65,3.48) {R37};
\node[note, anchor=west] at (10.95,3.48) {2 of 4};
\node[anchor=west, font=\fontsize{5.8}{6.8}\selectfont\bfseries] at (7.5,3.1) {set selection};
\node[chip] at (9.0,3.1) {R12}; \node[chip] at (9.55,3.1) {R37}; \node[chip] at (10.1,3.1) {R58}; \node[chip] at (10.65,3.1) {R61};
\node[note, anchor=west, text=cStdD] at (10.95,3.1) {all 4};
% ---------------------------------------------------------------- (4) widened path
\draw[panel, draw=cWidD, fill=cWid] (7.35,0) rectangle (11.75,2.45);
\node[num] at (7.6,2.18) {4}; \node[title] at (7.77,2.18) {Wider evidence access};
\draw[rounded corners=2pt, draw=cWidD!70, fill=white] (7.5,0.98) rectangle (11.6,1.9);
\node[body=3.9cm] at (7.6,1.8) {\textbf{Lightweight reader.} wider index lines (600 characters) and one explicit note per member of the set.};
\draw[rounded corners=2pt, draw=cWidD!70, fill=white] (7.5,0.15) rectangle (11.6,0.83);
\node[body=3.9cm] at (7.6,0.73) {\textbf{Strong reader.} the selected rounds plus the 40 most similar rounds.};
\draw[arr] (D.east) -- (7.15,1.25) |- (7.35,4.1);
\node[font=\fontsize{5.6}{6}\selectfont, anchor=south] at (6.78,1.27) {yes};
\draw[arr] (D.south) |- (7.35,0.62);
\node[font=\fontsize{5.6}{6}\selectfont, anchor=south] at (6.4,0.63) {no};
% ---------------------------------------------------------------- (5) reader
\draw[panel, draw=cReadD, fill=cRead] (12.15,0) rectangle (16,5.4);
\node[num] at (12.4,5.13) {5}; \node[title] at (12.57,5.13) {Read raw evidence};
\node[body=3.6cm] at (12.3,4.85) {Opens the selected rounds and answers from their original text and images.};
\foreach \x/\lab in {12.42/R12, 13.32/R37, 14.22/R58, 15.12/R61} {
  \draw[draw=black!40, fill=white, rounded corners=1pt] (\x,3.32) rectangle (\x+0.72,3.98);
  \fill[blue!15] (\x+0.06,3.38) rectangle (\x+0.66,3.82);
  \fill[blue!45] (\x+0.08,3.4) -- (\x+0.27,3.64) -- (\x+0.38,3.52) -- (\x+0.63,3.78) -- (\x+0.63,3.4) -- cycle;
  \node[font=\ttfamily\fontsize{5.2}{6}\selectfont, text=black!70] at (\x+0.36,3.17) {\lab};
}
\draw[rounded corners=2pt, draw=cReadD, fill=white, line width=0.6pt] (12.32,0.18) rectangle (15.83,2.45);
\node[align=left, text width=3.3cm, inner sep=0pt] at (14.07,1.32) {\textbf{Answer} combines the hotel invoice, the restaurant bill, the receipt on the D9 summary card, and the taxi receipt.};
\draw[arr] (14.07,2.98) -- (14.07,2.5);
\draw[arr] (11.75,4.1) -- (12.15,4.1);
\draw[arr] (11.75,1.25) -- (12.15,1.25);
\end{tikzpicture}
}
\caption{\textbf{Overview of \vgs{}.} A visibility gate, computed when the index is written, chooses between global evidence navigation and wider access adapted to the reader. Both paths read original rounds and images. Gold markers identify relevant receipts in the illustrative history.}
\label{fig:overview}
\end{figure*}
